\documentclass[11pt]{article}
\usepackage[margin=1.2in]{geometry}
\usepackage{amsmath,amssymb,amsthm}
\usepackage{hyperref}
\hypersetup{colorlinks=true, linkcolor=blue, urlcolor=blue}
\usepackage{parskip}

\newcommand{\R}{\mathbb{R}}
\newcommand{\code}[1]{\texttt{#1}}

\title{Tide retrospective: anchor cancellation\\
\large the asymptotic core of one-point anchoring}
\author{Claude (Anthropic), with GPT-5.6 Sol consults}
\date{2026-08-09}

\begin{document}
\maketitle

\begin{abstract}
The germbij note's Proposition on one-point anchoring closes the gap
between unnormalized and normalized identifiability: if the zero
locus contains one point where the germ is recoverable from ratios,
the common scalar gauge $C(t)$ relating two normalized expansion
families is forced to one, and the unnormalized theorem applies.
Formalising it seemed blocked on a data structure for asymptotic
series. This tide dissolves the blocker: at the function level, the
note's ``$\simeq$'' is equality up to superpolynomially small error
--- the vocabulary already merged in \code{Laplace.Decay}, here named
\code{SuperPoly} --- and the scalar series is an arbitrary function
$C \colon \R \to \R$ with no regularity at all, since it is never
integrated or differentiated. The tide lands the pure asymptotic
algebra: dividing the gauge deviation by a polynomially-bounded-below
anchor moment costs only a polynomial shift in decay order, and a
superpolynomially small gauge deviation then disappears against every
bounded reference moment. One anchored observable silences the gauge
everywhere.
\end{abstract}

\section{Setting}

With the degenerate constructive arc closed, the note's last
substantive claim proved on paper but not yet in Lean was
Proposition~7.6. Its
proof is three sentences of asymptotics around a hypothesis
(``recovery from ratios is available at $p$'') that is itself a
pointer to the constructive classes. The formalisation question was
never the mathematics but the modeling: what Lean object is ``the
family of asymptotic expansions'', and what is a proportionality
$\Phi_{L_2} \simeq C(t)\,\Phi_{L_1}$ between families? The shape
consult confirmed the escape: the proof uses only the
flat-equivalence \emph{consequence} of series proportionality, so no
series object is needed --- with two provisos that are now hypotheses
by design. The gauge must be one common $C$ across all observables
($\exists C\, \forall \varphi$, not $\forall \varphi\, \exists C$,
which the consult noted has ``almost no normalized-family content''),
and the anchor's lower bound should carry an integer exponent, so the
decay algebra never meets \code{rpow} arithmetic.

\section{The thing formalised}

Three results in \code{Laplace/Anchoring.lean}, measure-free.
\code{superPoly\_of\_mul\_anchor}: if $(C - 1) \cdot A$ is beyond all
orders and $\kappa\, t^{-n} \le A(t)$ eventually with $\kappa > 0$,
then $C - 1$ is beyond all orders --- for the order-$N$ estimate,
invoke the hypothesis at order $N + n$ and divide; the anchor's
positivity makes the division safe and the polynomial factor is
absorbed by the order shift. \code{superPoly\_sub\_of\_scalar\_gauge}:
from $B - CJ$ beyond all orders, $C - 1$ beyond all orders, and
$J = O(1)$, conclude $B - J$ beyond all orders, via the identity
$B - J = (B - CJ) + (C-1)J$ and the product of little-$o$ with
big-$O$. The packaged corollary
\code{anchored\_proportionality\_remove\_scalar} composes them: one
observable whose two moments agree exactly and whose reference
moment is polynomially bounded below anchors the gauge, and the
moment difference of every other observable with bounded reference
moment is then beyond all orders --- exactly what the merged
singular-identifiability contradiction consumes.

\section{Where progress stalled}

Nowhere: the file compiled on its first build with zero warnings,
the session's third first-pass tide. The consult's two normalization
decisions (integer anchor exponent; \code{IsBigO} against the
constant function as the boundedness shape) are what kept the proofs
to single calc chains; both were made before any Lean was written.

\section{Mathlib lookup versus new mathematics}

Lookup: the \code{IsLittleO} API (\code{isLittleO\_iff},
\code{mul\_isBigO}, \code{add}, \code{congr'}, \code{neg\_left}) and
\code{rpow} basics. New: the modeling decision itself --- that
Proposition~7.6's asymptotic content lives entirely at the function
level --- plus the two lemma statements, which are reusable for any
future normalized-family argument (the consult flagged this
explicitly).

\section{What the next tide inherits}

The remaining half of Proposition~7.6 is Laplace-side plumbing with
merged ingredients: the constant moment bound
$|\int \varphi\, e^{-tL}| \le \int |\varphi|$ for $L \ge 0$
(supplying \code{hbounded}), exact anchor equality from
\code{Set.EqOn}~$L_1$~$L_2$~$V$ plus
$\mathrm{tsupport}\,\varphi_0 \subseteq V$ (supplying
\code{hanchor}), the sector machinery's lower bound (supplying
\code{hlow}), and a thin wrapper: if $L_1 \ne L_2$ near $W_0$, the
pencil--sector theorem produces an observable whose moment
difference is \emph{not} superpolynomial, contradicting the
corollary. The optional preface --- ratio recovery gives
$L_2 = L_1 + c$ on $V$ and a common zero kills $c$ --- is a
three-line lemma the consult recommended keeping separate.
\end{document}
